\subsection{معماری و پیاده‌سازی}
پیاده‌سازی ماژول در فایل \lr{Codes/framework/compression\_engine/nli\_softmax.py} تعبیه شده و در هر دو رجیستری مکانیزم توجه و موتور فشرده‌سازی سیستم به ثبت رسیده است. اجزای کلیدی معماری عبارتند از:
\begin{itemize}
    \item \textbf{تابع سفارشی \lr{NLISoftmaxFunction}:} پیاده‌سازی \lr{torch.autograd.Function} با گذر پیشرو استاندارد و گذر پسرو بهینه‌شده به همراه مهار عددی گرادیان‌ها.
    \item \textbf{ماژول \lr{NLISoftmaxAttention}:} جایگزین ساختاری بلوک‌های توجه خودکار برای استقرار مستقیم بر روی لایه‌های انکودری و دیکودری.
    \item \textbf{موتور فشرده‌سازی \lr{NLISoftmaxCompressor}:} جایگزینی خودکار نمونه‌های \lr{nn.Softmax} در معماری ترانسفورمر، اعتبارسنجی هم‌ارزی پیشرو (\lr{Equivalence Validation}) با آستانه خطای \lr{max\_error} (پیش‌فرض $0.20$)، و مدیریت هوشمند لایه‌های استثنا نظیر سر طبقه‌بند و توابع هزینه.
\end{itemize}

\subsection{محدودیت‌ها و تحلیل شکست}
تحلیل ریشه‌ای شکست مدل در بنچ‌مارک \lr{RoBERTa} به عوامل زیر بازمی‌گردد:
\begin{itemize}
    \item \textbf{ناسازگاری ماسک‌گذاری توجه در ساختار تو در توی \lr{RobertaSelfAttention}:} معماری \lr{RoBERTa} از ماسک‌های باینری با مقادیر بسیار منفی بزرگ (نظیر $-10000.0$) برای نشانه‌های پدینگ استفاده می‌کند. جایگزینی مستقیم تابع بیشینه‌ساز نرم در زیرماژول‌های تو در تو، در ترکیب با فرآیند استخراج برچسب طبقه‌بند بدون هماهنگی سر خروجی، منجر به عدم محاسبه صحیح برچسب در حلقه ارزیابی دقت و گزارش دقت \lr{0.00\%} با وجود محاسبه شاخص F1 شده است.
    \item \textbf{پیچیدگی تزریق لایه‌ای در معماری‌های با مکانیزم توجه ادغام‌شده:} در مدل‌هایی که تابع بیشینه‌ساز نرم به عنوان عملگر درونی در متد \lr{forward} پیاده‌سازی شده و فاقد ماژول مجزای \lr{nn.Softmax} هستند، فشرده‌ساز ناگزیر به جایگزینی کل بلوک توجه بوده که این امر در صورت کوچکترین عدم انطباق تانسورهای خروجی با خطای ساختاری همراه می‌شود.
\end{itemize}
